% Generado a partir de calculus/Final_Bien/evidencia_P8.m (simulador corregido).
\subsection{Controlador con atascamiento-deslizamiento sobre el modelo lineal}

Para separar el efecto de las no linealidades no suaves del de la no linealidad de la fuerza electromagnética, el PII se probó sobre el modelo linealizado en $x_1^*=-2{,}5$~cm al que se añadieron la zona muerta y la fricción de Karnopp, con la misma prueba que se emplea en el capítulo~\ref{chap:resultados} para la regulación del modelo no lineal ($-2\to-3$~cm, $u\in[0;\,3{,}5]$~V). La figura~\ref{fig:p8_lineal_ad} compara ambos resultados. El comportamiento es prácticamente el mismo: sobrepaso de $24{,}4$\,\% frente a $23{,}6$\,\%, entrada en la banda de $\pm5$\,\% en $0{,}41$ frente a $0{,}40$~s y un ciclo límite de $0{,}028$ frente a $0{,}033$~cm pico a pico. La única diferencia apreciable es el voltaje de equilibrio final ($2{,}36$~V en el modelo lineal frente a $2{,}39$~V), consecuencia del error de linealización a $0{,}5$~cm del punto de operación. Esto indica que el ciclo límite se debe a la zona muerta y a la fricción estática, y no a la no linealidad de la fuerza electromagnética.

\begin{figure}[ht]
    \centering
    % Estilo común de las figuras de evidencia P8 (pgfplots).
\pgfplotsset{
  pocho/.style={
    width=0.92\textwidth, height=5.2cm,
    grid=major, grid style={gray!25},
    tick label style={font=\small}, label style={font=\small}, legend style={font=\small},
    /pgf/number format/use comma,
    scaled ticks=false,
    y tick label style={/pgf/number format/fixed},
    table/col sep=comma, unbounded coords=jump,
    every axis plot/.append style={line join=round},
  },
}
\definecolor{tray}{RGB}{31,94,166}
\definecolor{refc}{RGB}{40,40,40}
\definecolor{ctl2}{RGB}{196,96,40}
\definecolor{ver}{RGB}{46,133,64}

\begin{tikzpicture}
\begin{axis}[pocho, name=principal, xmin=0, xmax=40, ymin=-3.35, ymax=-1.85, xlabel={},
    ylabel={posición [\si{cm}]}, legend style={at={(0.99,0.97)}, anchor=north east}, legend cell align=left]
  \addplot[ctl2, line width=0.7pt] table[x=t, y=x_nl] {images/p8_results/tikz/lineal_ad.csv}; \addlegendentry{no lineal con AD}
  \addplot[tray, line width=1pt] table[x=t, y=x_lin] {images/p8_results/tikz/lineal_ad.csv}; \addlegendentry{lineal con AD}
\end{axis}
\begin{axis}[pocho, at={(principal.south west)}, anchor=north west, yshift=-0.45cm, height=4.4cm,
    xmin=0, xmax=40, ymin=-0.2, ymax=3.8, xlabel={tiempo [\si{s}]}, ylabel={$u$ [\si{V}]}]
  \addplot[ctl2, line width=0.7pt] table[x=t, y=u_nl] {images/p8_results/tikz/lineal_ad.csv};
  \addplot[tray, line width=0.7pt] table[x=t, y=u_lin] {images/p8_results/tikz/lineal_ad.csv};
\end{axis}
\end{tikzpicture}

    \caption[Modelo lineal con atascamiento-deslizamiento]{Regulación con zona muerta y fricción de Karnopp sobre el modelo lineal (en $x_1^*=-2{,}5$~cm) y sobre el modelo no lineal. Arriba, posición; abajo, señal de control.}
    \label{fig:p8_lineal_ad}
\end{figure}

\FloatBarrier
